\documentclass[11pt]{article}
\usepackage[a4paper,margin=25mm]{geometry}
\usepackage{amsmath,amssymb,mathtools}
\usepackage{fontspec,unicode-math}
\setmainfont{Latin Modern Roman}
\setsansfont{Latin Modern Sans}
\setmathfont{Latin Modern Math}
\usepackage{xurl,hyperref}
\hypersetup{colorlinks=true,linkcolor=blue,urlcolor=blue,pdftitle={Problem statement}}
\newcommand{\sref}[2]{\hyperlink{source:#1:#2}{[S#2]}}
\newcommand{\cataloguescope}[1]{\paragraph{Mathematical question.}#1}
\begin{document}
\section{Explicit full-dimensional Julia set}\label{statement}

\subsection{Definitions and mathematical statement}

\paragraph{Shared notation.}$\mathbb N=\{1,2,\ldots\}$, and $\mathbb Z,\mathbb Q,\mathbb R,\mathbb C$ denote the integers, rationals, reals and complex numbers. Any other conventions are specified locally.


For a rational map $f:\widehat{\mathbb C}\to\widehat{\mathbb C}$ of degree at least two, $J(f)$ is the complement of the largest open set on which the iterates form a normal family. Use the spherical metric for Hausdorff dimension and the spherical area measure for Lebesgue measure. The Hausdorff dimension $\dim_H E$ is the threshold exponent for Hausdorff measures $\mathcal H^s(E)$, defined using covers by sets of small diameter and sums of their $s$th powers. Planar Lebesgue measure is denoted by $m_2$.
\paragraph{Statement recorded in the supplied source.}
Find an explicit rational map $f$ with $\dim_HJ(f)=2$. When a Julia set has Hausdorff dimension two but Lebesgue measure zero, what is a natural geometric measure on it? \sref{statement}{2}
\subsection{Short English statement}
Give an explicit rational Julia set of dimension two, and identify a natural geometric measure for such sets when their ordinary area is zero.
\subsection{Sources}
\begin{description}
\item[\hypertarget{source:statement:1}{S1}] ULAM.AI and the UnsolvedMath Contributors, UnsolvedMath: A Curated Collection of Open Mathematics Problems, catalogue entry. CC BY 4.0. \url{https://huggingface.co/datasets/ulamai/UnsolvedMath}.
\item[\hypertarget{source:statement:2}{S2}] Mikhail Lyubich, ``Measure and Dimension of Julia Sets'', in Ben Bielefeld and Mikhail Lyubich (eds.), Problems in Holomorphic Dynamics (1992), Problem 5, pp. 27; arXiv:math/9205209v1. \url{https://arxiv.org/pdf/math/9205209}.
\end{description}
\label{end:statement}
\end{document}
